Name some essential properties of prime ideals:


A proper ideal \( P \) in a commutative ring \( R \) with unity is called a **prime ideal** if whenever \( a, b \in R \) such that \( ab \in P \), then either \( a \in P \) or \( b \in P \).
\(
P \text{ is prime } \iff ab \in P \Rightarrow a \in P \text{ or } b \in P
\)

In the ring of integers \(\mathbb{Z}\), the prime ideals are of the form \( (p) \) where \( p \) is a prime number.
\(
P \subseteq \mathbb{Z} \text{ is prime } \iff P = (p) \text{ for some prime } p
\)

Prime Ideal in a Polynomial Ring
If \( P \) is a prime ideal in a commutative ring \( R \), then \( P[x] \) is a prime ideal in \( R[x] \).
\(
P \text{ is prime in } R \Rightarrow P[x] \text{ is prime in } R[x]
\)

Localization
If \( P \) is a prime ideal in a commutative ring \( R \), then the localization \( R_P \) has a unique maximal ideal, which is \( PR_P \).
\(
P \text{ is prime in } R \Rightarrow PR_P \text{ is the unique maximal ideal of } R_P
\)

Intersection of Prime Ideals
The intersection of all prime ideals containing a given ideal \( I \) in a ring \( R \) is called the radical of \( I \), denoted \( \sqrt{I} \).
\(
\sqrt{I} = \bigcap_{\substack{P \supseteq I \\ P \text{ prime}}} P
\)

Correspondence Theorem
There is a one-to-one correspondence between the prime ideals of \( R/I \) and the prime ideals of \( R \) that contain \( I \).
\(
P \text{ is prime in } R/I \iff P \text{ is prime in } R \text{ and } I \subseteq P
\)

Prime Avoidance Lemma
If \( P_1, P_2, \ldots, P_n \) are prime ideals in a ring \( R \), and \( I \) is an ideal such that \( I \subseteq \bigcup_{i=1}^n P_i \), then \( I \subseteq P_i \) for some \( i \).
\(
I \subseteq \bigcup_{i=1}^n P_i \Rightarrow I \subseteq P_i \text{ for some } i
\)

Nilradical and Prime Ideals
The nilradical of a ring \( R \), denoted by \( \sqrt{(0)} \), is the intersection of all prime ideals in \( R \).
\(
\sqrt{(0)} = \bigcap_{\substack{P \text{ prime}}} P
\)

Quotient Rings
If \( P \) is a prime ideal in \( R \), then the quotient ring \( R/P \) is an integral domain.
\(
P \text{ is prime in } R \Rightarrow R/P \text{ is an integral domain}
\)

Extension and Contraction
Let \( R \to S \) be a ring homomorphism. If \( P \) is a prime ideal in \( R \), then its extension \( P^e \) in \( S \) is also a prime ideal if the homomorphism is injective.
\(
P \text{ is prime in } R \text{ and } R \to S \text{ is injective} \Rightarrow P^e \text{ is prime in } S
\)


Give some examples of prime ideals in selected rings:

Prime Ideals in \(\mathbb{Z}\)
In the ring of integers \(\mathbb{Z}\), the prime ideals are of the form \( (p) \) where \( p \) is a prime number. For example:
\(
(2), (3), (5), (7), \ldots
\)

Prime Ideals in Polynomial Rings \(\mathbb{Z}[x]\)
- The ideal \((2)\) in \(\mathbb{Z}[x]\) is a prime ideal.
- The ideal \((x)\) in \(\mathbb{Z}[x]\) is also a prime ideal because \(\mathbb{Z}[x]/(x) \cong \mathbb{Z}\), which is an integral domain.

Prime Ideals in Polynomial Rings over Fields \(\mathbb{K}[x]\)
- If \(\mathbb{K}\) is a field, then the ideal \((x-a)\) in \(\mathbb{K}[x]\) for any \(a \in \mathbb{K}\) is a prime ideal because \(\mathbb{K}[x]/(x-a) \cong \mathbb{K}\), which is a field and hence an integral domain.
- The ideal \((x^2 + 1)\) in \(\mathbb{R}[x]\) is a prime ideal because \(\mathbb{R}[x]/(x^2 + 1) \cong \mathbb{C}\), which is a field.

Prime Ideals in Rings of Integers in Number Fields
- In the ring of Gaussian integers \(\mathbb{Z}[i]\), the ideal \((1+i)\) is a prime ideal because \(\mathbb{Z}[i]/(1+i) \cong \mathbb{F}_2\), which is a field.
- In the ring of Eisenstein integers \(\mathbb{Z}[\omega]\), where \(\omega = e^{2\pi i / 3}\), the ideal \((1 - \omega)\) is a prime ideal.

Prime Ideals in \(\mathbb{Z}/6\mathbb{Z}\)
- The ideal \((2)\) in \(\mathbb{Z}/6\mathbb{Z}\) is a prime ideal because \((\mathbb{Z}/6\mathbb{Z})/(2) \cong \mathbb{Z}/3\mathbb{Z}\), which is a field.
- Similarly, the ideal \((3)\) in \(\mathbb{Z}/6\mathbb{Z}\) is a prime ideal because \((\mathbb{Z}/6\mathbb{Z})/(3) \cong \mathbb{Z}/2\mathbb{Z}\), which is a field.

Prime Ideals in Noncommutative Rings
- In the ring of \(n \times n\) matrices over a field \(\mathbb{K}\), \(M_n(\mathbb{K})\), the only prime ideal is \((0)\) when \(n > 1\).

These examples illustrate various contexts in which prime ideals appear, showcasing their importance in different areas of algebra.


Give some examples of 
Intersection of Ideals \(I \cap J\)
- The intersection of two ideals \(I\) and \(J\) in a ring \(R\) is always an ideal of \(R\).
\(
I \cap J = \{ r \in R \mid r \in I \text{ and } r \in J \}
\)

Sum of Ideals \(I + J\)
- The sum of two ideals \(I\) and \(J\) in a ring \(R\) is always an ideal of \(R\).
\(
I + J = \{ r \in R \mid r = i + j \text{ for some } i \in I \text{ and } j \in J \}
\)

Product of Ideals \(IJ\)
- The product of two ideals \(I\) and \(J\) in a ring \(R\) is always an ideal of \(R\).
\(
IJ = \left\{ \sum_{k=1}^n a_k b_k \mid a_k \in I, b_k \in J, n \in \mathbb{N} \right\}
\)

Quotient of Ideals
- If \(I \subseteq J\), then the quotient \(J/I\) forms an ideal in the quotient ring \(R/I\).

Radical of an Ideal \(\sqrt{I}\)
- The radical of an ideal \(I\) in a ring \(R\) is an ideal of \(R\).
\(
\sqrt{I} = \{ r \in R \mid r^n \in I \text{ for some } n \in \mathbb{N} \}
\)

Extended and Contracted Ideals
- Let \(f: R \to S\) be a ring homomorphism. If \(I\) is an ideal of \(R\), the extended ideal in \(S\) is defined as:
 \(
 I^e = f(I)S = \{ s \in S \mid s = f(i) \cdot s' \text{ for some } i \in I \text{ and } s' \in S \}
 \)
 - The contraction of an ideal \(J\) in \(S\) is defined as:
 \(
 J^c = f^{-1}(J) = \{ r \in R \mid f(r) \in J \}
 \)
 Both the extension and contraction are ideals in \(S\) and \(R\), respectively.

Union of Ideals \(I \cup J\)
- The union of two ideals \(I\) and \(J\) is not necessarily an ideal unless one ideal is contained in the other, i.e., \(I \subseteq J\) or \(J \subseteq I\).
\(
I \cup J \text{ is an ideal if and only if } I \subseteq J \text{ or } J \subseteq I
\)


Describe and relate principal ideal rings, unique factorization domains, euclidean domains.

Definitions

Principal Ideal Domain (PID)
A ring \( R \) is a principal ideal domain if every ideal in \( R \) is principal, i.e., can be generated by a single element.
\( 
R \text{ is a PID} \iff \forall \text{ ideal } I \subseteq R, \exists a \in R \text{ such that } I = (a)
\) 

Unique Factorization Domain (UFD)
A ring \( R \) is a unique factorization domain if every element can be factored into irreducibles uniquely, up to units and order.
\( 
R \text{ is a UFD} \iff \forall a \in R \setminus \{0\}, \exists \text{ irreducibles } p_1, p_2, \ldots, p_n \in R \text{ such that } a = u p_1 p_2 \cdots p_n
\) 
where \( u \) is a unit, and the factorization is unique up to order and units.

Euclidean Domain (ED)
A ring \( R \) is a Euclidean domain if there exists a Euclidean function \( \phi: R \setminus \{0\} \to \mathbb{N} \) such that for any \( a, b \in R \) with \( b \neq 0 \), there exist \( q, r \in R \) such that:
\( 
a = bq + r \quad \text{with either} \quad r = 0 \quad \text{or} \quad \phi(r) < \phi(b)
\) 

All of them are noetherian (this property is lost 
by GCD domains which contains the previous categories)

Examples
- **Euclidean Domain**: \(\mathbb{Z}\), \( \mathbb{F}[x] \) where \( \mathbb{F} \) is a field.
- **Principal Ideal Domain**: \(\mathbb{Z}\), \( \mathbb{F}[x] \), but not \(\mathbb{Z}[x]\).
- **Unique Factorization Domain**: \(\mathbb{Z}\), \( \mathbb{F}[x] \), \(\mathbb{Z}[x]\), , where \( \mathbb{F} \) is a field.

### Summary Diagram
\( \text{Euclidean Domain (ED)} \subseteq \text{Principal Ideal Domain (PID)} \subseteq \text{Unique Factorization Domain (UFD)}\) 


Define the prerequisites of an affine/projective algebraic variety and itself.


For an algebraically closed field \(K\) and a natural number \(n\), let \(A^n\) be an affine \(n\)-space over \(K\), identified to 
\(K^{n}\). For each set \(S\) of polynomials in \(K[x_1, \dots, x_n]\), define the zero-locus \(Z(S)= \{ x\in \mathbf{A} ^{n}\mid f(x)=0{\text{ for all }}f\in S\}\).
A subset \(V\) of \(A^n\) is called an affine algebraic set if \(V = Z(S)\) for some \(S\).
A nonempty affine algebraic set \(V\) is called irreducible, if it cannot be written as the union of two proper algebraic subsets.
An irreducible affine algebraic set is also called an affine variety.

For a projective variety, the underlying space is \(A^{n+1}\) after identifying points \(x \sim x' \iff x = \lambda x, \lambda \in K\), for an arbitrary \(f \in K[x_1, \dots, x_n]\) this creates an issue since there can be such equivalent points \(x \sim x'\) with \(f(x) = 0\), but \(f(x') \neq 0\) (e.g. \(K = \mathbb{R}, f(x) = (x - 2)^2\), \(f(1) = 1, f(2 \cdot 1) = 0\), this can be remedied by enforcing \(f\) to be homogenous, hence \(S\) may only contain such polynomials. Ceteris paribus \(V = Z(S)\) irreducible, is a projective variety.



the only issue arises that \(f = \lambda f, \forall \lambda \in K\), therefore we need to be able to factor out arbitrary constant
out of the polynomial, which requires them to be homogenous



Define the Zariski topology for affine 


through the choice of an affine coordinate system. 

Let \(f \in K[x_1, \dots, x_n]\) can be viewed as \(K\)-valued functions on \(A^{n}\) by evaluating \(f\) at the points in \(A^n\), i.e. by choosing values in \(K\) for each \(x_i\). 
\(Z(S)\) to be the set of points in \(A^n\) on which the functions in \(S\) simultaneously vanish, that is to say



State and proof the existence and uniqueness of the tensor product of modules:

Let \( R \) be a commutative ring and \( M \) and \( N \) be \( R \)-modules. The tensor product of \( M \) and \( N \), denoted \( M \otimes_R N \), is an \( R \)-module equipped with a bilinear map \( \otimes: M \times N \to M \otimes_R N \) satisfying the following universal property:

For any \( R \)-module \( P \) and any bilinear map \( f: M \times N \to P \) (i.e., a map that is \( R \)-linear in each argument), there exists a unique \( R \)-linear map \( \bar{f}: M \otimes_R N \to P \) such that the following diagram commutes:
\[
\begin{array}{ccc}
M \times N & \xrightarrow{\otimes} & M \otimes_R N \\
 & \searrow f & \downarrow \bar{f} \\
 & & P
\end{array}
\]
This means \( f = \bar{f} \circ \otimes \).

To show the existence of the tensor product, we construct it explicitly. Consider the free \( R \)-module \( F \) generated by the set \( M \times N \). Let \( F \) be the free \( R \)-module with basis elements \( [m, n] \) for \( m \in M \) and \( n \in N \). Define a submodule \( K \) of \( F \) generated by the elements of the following forms:
1. \( [m_1 + m_2, n] - [m_1, n] - [m_2, n] \) for all \( m_1, m_2 \in M \) and \( n \in N \),
2. \( [m, n_1 + n_2] - [m, n_1] - [m, n_2] \) for all \( m \in M \) and \( n_1, n_2 \in N \),
3. \( [rm, n] - r[m, n] \) for all \( r \in R \), \( m \in M \), and \( n \in N \),
4. \( [m, rn] - r[m, n] \) for all \( r \in R \), \( m \in M \), and \( n \in N \).

Define \( M \otimes_R N \) as the quotient module \( F / K \). Let \( \otimes: M \times N \to M \otimes_R N \) be the map sending \( (m, n) \) to the equivalence class \( [m, n] + K \). 

### Verification of the Universal Property

Let \( P \) be any \( R \)-module, and let \( f: M \times N \to P \) be a bilinear map. We need to show that there is a unique \( R \)-linear map \( \bar{f}: M \otimes_R N \to P \) such that \( f = \bar{f} \circ \otimes \).

Define \( \bar{f}: F \to P \) on the basis elements by \( \bar{f}([m, n]) = f(m, n) \). This map \( \bar{f} \) is \( R \)-linear because \( f \) is bilinear. We can straightforward verify that \( \bar{f} \) annihilates the submodule \( K \) by checking it for the generators of \( K \).

Hence it induces a well-defined map \( \bar{f}: M \otimes_R N \to P \). By construction, \( \bar{f}([m, n]) = f(m, n) \), so \( f = \bar{f} \circ \otimes \).

For the uniqueness,
suppose \( M \otimes'_R N \) is another \( R \)-module with a bilinear map \( \otimes': M \times N \to M \otimes'_R N \) satisfying the universal property. There exists a unique \( R \)-linear map \( \phi: M \otimes_R N \to M \otimes'_R N \) such that \( \phi \circ \otimes = \otimes' \).

Similarly, there exists a unique \( R \)-linear map \( \psi: M \otimes'_R N \to M \otimes_R N \) such that \( \psi \circ \otimes' = \otimes \).

Consider \( \phi \circ \psi: M \otimes'_R N \to M \otimes'_R N \). Since \( \phi \circ \psi \circ \otimes' = \phi \circ \otimes = \otimes' \), and by the uniqueness part of the universal property, \( \phi \circ \psi \) must be the identity on \( M \otimes'_R N \). Similarly, \( \psi \circ \phi \) must be the identity on \( M \otimes_R N \). Thus, \( \phi \) and \( \psi \) are isomorphisms, proving the uniqueness up to isomorphism of \( M \otimes_R N \).



Prove some natural isomorphisms of the tensor product, like the associative, symmetry and direct product ones. 

For \(A\)-modules \(M, N\) there is a natural isomorphism
\(M \otimes N \simeq N \otimes\), which is \(\phi(m \otimes n) = n \otimes m\) (note that this does not imply commutativity of the elementary tensors).
Proof:
\(\kappa : M \times N \to N \otimes M, \kappa((m, n)) = n \otimes m \) is bilinear, by the universal property \(\phi : M \otimes N \to N \otimes M\) must then exist, similarly \(\kappa^{\text{op}}(n, m) = m \otimes n\) exist as well and the induced \(\phi^{\text{op}}\) is readily seen to be invers to \(\phi\).

For \(A\)-modules \(L, M, N\) there is a natural isomorphism
\(L \otimes (M \otimes N) \simeq (L \otimes M) \otimes N\), which is \(\phi(l \otimes (m \otimes n)) = (l \otimes m) \otimes n\).
Proof:
\(f_s : M \times N \to (L \otimes M) \otimes N, f_s(m, n) = (s \otimes m) \otimes n\) is bilinear, so by the universal property there is some
\(\tilde{f}_s : (M \otimes N) \to (L \otimes M) \otimes N\) matching \(f_s\) and further this can be extended to some \(\tilde{f} : L \times (M \otimes N) \to (L \otimes M) \otimes N\), which becomes the isomorphism from \(L \otimes (M \otimes N)\) in question via the universal property.

Let \(f : M \to M'\), \(g : N \to N'\) be \(A\)-linear maps, then there exists a unique map \(h : M \otimes N \to M' \otimes N'\) with \(h(m \otimes n) = f(m) \otimes g(n)\).
Clearly \(\bar{h}\) is linear (here the domain is still \(M \times N\)), so there is a unique bilinear map \(h\) satisfying the required.

For the direct sum of modules \(\bigoplus N_i\) it holds \(\bigoplus_{i} M \otimes N_i \simeq M \otimes \bigoplus N_i\).

Consider the family of morphisms \(\psi_i : M \otimes N_i \to M \otimes \bigoplus N_i, \psi_i (m \otimes n_i) = m \otimes (0, \dots, n_i, \dots, 0)\), clearly this
family is linear so by the universal property of the direct sum there is a unique \(\theta : \bigoplus_{i} M \otimes N_i \to M \otimes \bigoplus N_i\), which restricts to each of the \(\psi_i\). Conversely \(\phi_i : M \otimes \bigoplus N_i \to M \otimes N_i, \phi_i(m \otimes (n_1, n_2, \dots, n_i, \dots)) = m \otimes n_i \) extends
by the universal property of the direct product (which holds also for direct sums) these also extends to some \(\gamma : M \otimes \bigoplus_i N_i \to \bigoplus M \otimes N_i\) and its readily seen that \(\theta^{-1} = \gamma\).



Prove and state the general theorem behind \(V^{*} \otimes V \simeq \text{Mat}_n(K)\).

There is a canonical homomorphism \(\mu : M* \otimes N \to \text{Hom}_A(M, N)\), where \(M^{*}\) is the dual modul.
In fact if \(M\) is free of finite rank, then \(\mu\) is an isomorphism.

Let the homomorphism be defined pointwise via \(\mu(m* \otimes n)(x) = m*(x) n\) then this is a linear map from \(M* \times N \to \text{Hom}_A(M, N)\) so by the universal
property the desired \(\mu\) exists. Additionally if \(M\) is free with basis \(\{m_j\}_{j \in [J]}\), \(M^{*}\) has a dual basis \(\{m^{*}_j\}_{j \in [J]}\), s.t.
\(m^{*}_j(m_i) = \delta_{ij}\). let \(\mu' : \text{Hom}_A(M, N) \to M* \otimes N\) be defined as \(\mu'(f) = \sum_{j \in [J]} m^{*}_j f(x_j)\), it follows that \(\mu'\)
is inverse to \(mu\).


How does the tensor product act on free modules?

Let \(M\) be a \(A\)-module (potentially not free) and \(N\) free with basis \(\{n_j\}_{j \in J}\)
then every element in \(M \otimes N\) has unique representation of the shape
\(\sum_{j \in J} m_j \otimes n_j\) with \(m_j = 0\) for nearly all \(j\).

If \(|J| = 1\) its easy to see that \(M \otimes N \simeq N\) and otherwise \(N = \bigoplus_{j \in J} A n_j\)
thus by the distributivity over direct sums the rest follows.

In particular if \(M\) is also free with basis \(\{m_i\}_{i \in I}\) then \(\{m_i \otimes n_j\}_{(i, j) \in I \times J}\) is a
basis of \(M \otimes N\), which a free \(A\)-module of rank \(|I||J|\).

Let \(f: M \times N \to P\) be an \(A\)-bilinear mapping. For each \(m \in M\) the mapping
\(n \to f(m, n)\) of \(N\) into \(P\) is \(A\)-linear, hence \(f\) gives rise to a mapping 
\(M \to \text{Hom}(N, P)\) which is \(A\)-linear because \(f\) is linear in the variable \(m\). Conversely
any \(A\)-homomorphism \(\phi: M \to \text{Hom}_A(N, P)\) defines a bilinear map, namely
\(\psi: M \times N \to P, (m, n) \mapsto \phi(x)(y)\). Hence the set \(S\) of A-bilinear mappings \(M \times N \to P\) is in
natural one-to-one correspondence with \(\text{Hom}(M, \text{Hom}(N, P))\). On the other
hand \(S\) is in one-to-one correspondence with \(\text{Hom}(M \otimes N, P)\), by the de-
fining property of the tensor product. Hence we have a canonical isomorphism
\(\text{Hom}(M \otimes N, P) \simeq \text{Hom}(M, \text{Hom}(N, P))\).

Given the exact sequence
\(M' \to M \to M'' \to 0\) with \(f : M' \to M\), \(g : M \to M''\) then the sequence \(M' \otimes N \to M \otimes N \to M'' \otimes N \to 0\) is exact
with modified morphisms \(f \otimes \text{id}\), \(g \otimes \text{id}\). 
Let \(E\) be the SES then \(\text{Hom}(E, \text{Hom}(N, P)\) is exact by the previous and via the isomorphism \(\text{Hom}(E \times N, P)\) is exact implying
\(E \times N\) is exact, where we used the arbitrariness of \(P\) (only if the sequence is exact for all \(P\) the left-argument of Hom is also exact).

In particular covariant \(\text{Hom}(X, -)\) is a right exact functor if and only if \(X\) is projective and contravariant \(\text{Hom}(-, X)\) left exact iff. \(X\) 
is injective.

Note that this does not imply the exactness of \(0 \to M' \to M \to M'' \to 0\) tensored with \(N\). However this does hold if the sequence is split.

Let \(f : R \to S\) be an R-algebra and M be an R-module. We consider the \(R\)-module
\( S \otimes_{R} M\), sometimes called \(f_{*}M\), and equip it with an \(S\)-module structure via
\(s (s' \otimes m) = ss' \otimes m\), which is the extension of scalars.

Similar for \(m\) an \(S\)-module one can construct the restriction of scalars.

This enjoys the following universal property:

Let \(S\) be an \(R\)-algebra via \(f : R \maps S\) and let \(N\) be an \(S\)-module and suppose there
is an \(R\)-linear map \(\phi : M \to N\). Then there is a unique \(S\)-linear map \(\Phi : S \otimes_{R} M \to N\) so that
\(\phi = \Phi \circ \iota\), where \(\iota : M \to S \otimes_{R} M\) is the natural inclusion.

If \(S\) is commutative, then we only need to require that \(f\) is a ring homomorphism which already implies \(f(R) \subset Z(S)\).
In particular if \(M\) is a free \(A\)-module then \(S \otimes B\) is again free with basis \(\{m_i \otimes 1_{S}\}_{i \in I}\).

A classical example is \(S \otimes A[x] \simeq S[x]\).
